\section{噪声调度：为什么中间区域最值得关注}

\subsection{为什么要重新采样 $t$}

训练 diffusion 或 rectified flow 时，不能把所有噪声水平一视同仁。最难的区域通常在中间噪声附近：信息既不够清晰，也没到完全随机。这个区域才是模型真正要“学会平均”的地方。

\begin{figure}[H]
\centering
\includegraphics[width=0.95\textwidth]{slides-images/slide-072.jpg}
\caption{Optimal prediction：网络在中间噪声上最难，因为同一个 $x_t$ 可能来自很多不同的 $(x,z)$ 对\protect\footnotemark}
\end{figure}
\footnotetext{Slides 第73页。}

\begin{figure}[H]
\centering
\includegraphics[width=0.95\textwidth]{slides-images/slide-075.jpg}
\caption{极端噪声很容易：$t=1$ 时更接近纯噪声，$t=0$ 时更接近真实数据\protect\footnotemark}
\end{figure}
\footnotetext{Slides 第76页。}

\begin{importantbox}{为什么中间噪声最难}
在中间噪声区域，输入还保留一部分结构，但又不足以唯一决定答案。模型在这里要处理的是“多对一”的平均问题，而不是简单复原。这也是噪声 schedule 设计最关键的地方。
\end{importantbox}

\subsection{非均匀 schedule 不是细节，而是模型设计}

如果训练时对所有噪声水平等权处理，就会让模型在大量“容易样本”上浪费算力。于是更好的做法是用非均匀噪声分布，比如 logit-normal sampling，甚至在高分辨率任务里把分布往更高噪声区域偏一点。

\begin{figure}[H]
\centering
\includegraphics[width=0.95\textwidth]{slides-images/slide-079.jpg}
\caption{噪声调度的结论：解决办法不是均匀采样，而是让训练更多覆盖中间噪声区域\protect\footnotemark}
\end{figure}
\footnotetext{Slides 第80页。}

\begin{figure}[H]
\centering
\includegraphics[width=0.95\textwidth]{slides-images/slide-081.jpg}
\caption{常见选择：logit-normal 采样，并在高分辨率任务里适当上移噪声分布\protect\footnotemark}
\end{figure}
\footnotetext{Slides 第82页。}

\begin{warningbox}{schedule 影响的不只是收敛速度}
schedule 不只是“调收敛快慢”，它直接影响网络看到的数据难度分布。换句话说，schedule 改变了你到底在教模型先学什么、后学什么。
\end{warningbox}

